\chapter{Conclusion and Future Work} \label{chapter:ch7_conclusion}

\section{Conclusion}

This thesis asked whether collaborative-robot tasks that change frequently can be authored more directly by specifying spatial intent in the real workspace instead of relying only on teach pendant, offline programming, or hand-guiding workflows. Within prototype scope, the answer is positive for selected task-oriented workflows. The developed system shows that poses and trajectories entered in the workspace can be combined with scene detection, task-specific interpretation, and robot execution to produce executable procedures on a real robot without exposing the operator to a full robot programming environment.

The main contribution is a reusable task-oriented workflow rather than only a tracked input device or two isolated demonstrations. Within a selected use case, the same workflow supports authoring, confirmation into executable commands, execution, visualization, and persistence. A shared runtime with explicit interfaces allows both pick-and-place and seam welding to reuse the same infrastructure.

The result nevertheless remains a proof of concept. The evaluation is qualitative, physical validation is limited, and the thesis does not include a formal user study, a baseline comparison against teach pendant or offline programming workflows, or a quantitative characterization of tracking accuracy, repeatability, or productivity. Within that boundary, seam welding is the clearest practical fit because the authored spatial input maps directly to a welding trajectory. Pick-and-place mainly serves as flexibility evidence because it shows that authored intent can be combined with refreshed scene information at execution time, but it depends more strongly on the narrow scene model used in the prototype.

\section{Near-Term Future Work}

A first near-term step is to strengthen sensing, tracking, and calibration. More robust pen tracking, better calibration diagnostics, fuller use of pen-side IMU data, and a richer scene estimate than the current fiducial-box representation would directly improve the present workflow. The same work should also better characterize tracking stability, scene uncertainty, and repeated scene acquisition in the real workspace.

A second near-term step is to complete task execution beyond the current prototype level. This includes a clearer abstraction for process tools, explicit motion validation and collision checking for authored tasks, and simulated execution for safer inspection before real-robot use. For tasks that do not depend on runtime scene adaptation, it is also worth exploring a path from authored commands to a lower-level robot-side export or compilation target.

The current prototype also needs broader validation and platform completion. Broader physical testing and bounded characterization of repeatability, tracking accuracy, and scene-estimation stability would make the present evidence easier to interpret. Cross-platform completion and a small amount of interface and orchestration cleanup belong in the same group because they strengthen the existing prototype rather than change the research direction.

\section{Longer-Term Research Directions}

Beyond prototype completion, one natural direction is richer input for spatial authoring. The architecture is not tied to the current pen-based interface, so alternative inputs such as hand tracking, VR or AR controllers, or combinations of speech and gesture could be explored within the same task-oriented model. This would help clarify which kinds of input are most suitable for different use cases and how much spatial intent can be specified naturally without losing control over the resulting robot behavior.

Another direction is to study how this authoring model fits existing robot-programming practice. One possibility is deeper integration with teach-pendant or controller workflows, so that spatial authoring becomes one programming tool among others rather than a separate frontend. Another is a formal user study focused on when this approach is actually useful, for which task classes it reduces setup friction, where its limits remain too restrictive, and how operators understand the boundary between selected use cases and more general robot programming.

The prototype should also be exercised on a broader range of task families and operating conditions. Extending the work toward richer welding scenarios, additional fabrication or surface-processing tasks, and more varied human-in-the-loop workflows would provide a stronger basis for judging that applicability boundary.

The present prototype provides a bounded starting point for that work, but broader claims still depend on the technical completion and evaluation outlined above.
